% =============================================================================
% Chapter 10: Results -- Part II
% =============================================================================
\mychapter{10. Results}

\mysection{10.1 Full-Volume Classical Thresholding Comparison}
Table~\ref{tab:podfam_classical} reports defect fraction, pore count, and mean equivalent pore
diameter for Otsu, Yen, and Bernsen thresholding, computed over the complete slice volume of every
core-validation PODFAM sample (749--900 slices).

\begin{table}[htbp]
\centering
\small
\caption{Full-volume classical thresholding comparison, all five core-validation PODFAM samples.
  sample\_03 and sample\_05 excluded from training (Sections~9.4, 9.5). Pixel units, see Section~12
  on voxel-size calibration.}
\label{tab:podfam_classical}
\begin{tabular}{lllll}
\toprule
\textbf{Sample} & \textbf{Method} & \textbf{Defect (\%)} & \textbf{Pore count} & \textbf{Mean $\varnothing$ (px)} \\
\midrule
sample\_01 (test)   & Otsu             & 0.003  & 997        & 5.36 \\
                     & Yen              & 42.535 & 1,743,429  & 5.15 \\
                     & Bernsen (label)  & 0.043  & 11,928     & 5.41 \\
sample\_02 (train)  & Otsu             & 2.993  & 336,468    & 8.20 \\
                     & Yen              & 38.132 & 443,339    & 10.19 \\
                     & Bernsen (label)  & 7.586  & 478,196    & 10.67 \\
sample\_03 (excl.)  & Otsu             & 0.067  & 8,850      & 7.87 \\
                     & Yen              & 45.319 & 1,737,239  & 5.19 \\
                     & Bernsen          & 0.230  & 51,467     & 5.55 \\
sample\_04 (val)    & Otsu             & 0.315  & 30,992     & 6.71 \\
                     & Yen              & 39.166 & 833,872    & 5.95 \\
                     & Bernsen (label)  & 1.994  & 103,238    & 9.00 \\
                     & U-Net (pred.)    & 3.239  & 122,587    & 11.28 \\
sample\_05 (excl.)  & Otsu             & 51.729 & 53,315     & 24.57 \\
                     & Yen              & 47.940 & 61,375     & 23.85 \\
                     & Bernsen          & 52.767 & 51,142     & 24.73 \\
\bottomrule
\end{tabular}
\end{table}

Yen produces a physically implausible defect fraction (38--45\%) on every sample without exception,
including the near-empty sample\_03, confirming, on an independent dataset, the over-segmentation
failure mode reported for Yen thresholding on AM XCT data~\citep{Kim2017}. Otsu and Bernsen disagree
with each other by one to two orders of magnitude on the same samples (e.g.\ sample\_02: Otsu 2.993\%
vs.\ Bernsen 7.586\%), a disagreement of the same order as the inter-method uncertainty separately
quantified in the literature~\citep{AppliedSciences2025}. The U-Net's predicted defect fraction on
its own validation sample, sample\_04 (3.239\%, Table~\ref{tab:podfam_classical}), is of the same
order as, though somewhat above, that sample's Bernsen label (1.994\%); as sample\_04 informed
checkpoint selection via early stopping, this number is reported for completeness rather than as an
independent generalisation result, that role is reserved for the fully held-out sample\_01
(Sections~10.3, 10.4).

\mysection{10.2 U-Net Training}
The detector was trained on sample\_02 (train), validated on sample\_04 (val), and tested on
sample\_01 (test), per the ranked assignment described in Section~9.4. Table~\ref{tab:podfam_epochs}
reports validation metrics at selected epochs; the best checkpoint (epoch 8) was used for all
subsequent evaluation. Figure~\ref{fig:training_curves} plots every one of the 18 logged epochs,
showing how the loss and each validation metric actually moved epoch to epoch, including fluctuation
that a selected-epoch summary table alone does not make visible.

\begin{table}[htbp]
\centering
\caption{2D U-Net training progression (dice\_focal loss, $\lambda = 0.5$), selected epochs. Full
  epoch-by-epoch log recorded via MLflow.}
\label{tab:podfam_epochs}
\begin{tabular}{llllll}
\toprule
\textbf{Epoch} & \textbf{Train loss} & \textbf{Val Dice} & \textbf{Val IoU} & \textbf{Val Recall} & \textbf{Note} \\
\midrule
1  & 0.346 & 0.497 & 0.353 & 0.996 & \\
2  & 0.289 & 0.521 & 0.372 & 0.891 & \\
4  & 0.215 & 0.585 & 0.434 & 0.874 & \\
8  & 0.115 & 0.600 & 0.452 & 0.954 & best checkpoint \\
18 & 0.058 & 0.559 & 0.407 & 0.938 & early-stopped (patience 10) \\
\bottomrule
\end{tabular}
\end{table}

\begin{figure}[htbp]
\centering
\includegraphics[width=0.95\linewidth]{training_curves_epoch_by_epoch.png}
\caption{Left: training and validation dice-focal loss per epoch. Right: validation Dice, IoU,
  recall, and training Dice per epoch. Dashed line marks epoch 8, the checkpoint selected for all
  downstream evaluation; training continued for a further 10 epochs under the patience-10
  early-stopping rule before stopping at epoch 18 without dice/IoU exceeding the epoch-8 value.}
\label{fig:training_curves}
\end{figure}

Training loss falls smoothly and monotonically across all 18 epochs (0.346 $\to$ 0.058), while
validation loss decreases far less smoothly and only until roughly epoch 8, after which it plateaus
and mildly oscillates (0.294 at epoch 8 vs.\ 0.254--0.306 for epochs 9--18), the classic
training/validation divergence pattern that early stopping is designed to catch. Validation Dice and
IoU follow the same shape: both rise unevenly through epoch 8 and then oscillate around a similar or
slightly lower level for the remaining 10 epochs (val Dice: 0.60 at epoch 8, ranging 0.51--0.59
afterward) rather than continuing to improve. Validation recall is the most volatile metric across
the whole run (oscillating between 0.857 and 0.996), consistent with a model still adjusting its
confidence threshold response on a validation set of a single volume, with only one sample in each of
the train/validation/test roles (Section~9.4), epoch-to-epoch metric swings reflect genuine sampling
variance in which patches were drawn each epoch, not measurement noise. Training Dice, by contrast,
climbs smoothly and near-monotonically throughout (0.783 $\to$ 0.921), the expected signature of a
model continuing to fit its training data after the validation metric it is being selected on has
stopped improving.

Against the thesis's stated acceptance criteria, validation recall (0.954) exceeds the 0.80 target;
validation Dice (0.600) and IoU (0.452) remain below the 0.75 and 0.60 targets respectively at the
point training was stopped.

\mysection{10.3 Held-Out Test Performance}
On the fully held-out sample\_01, the trained detector achieved Dice 0.147, IoU 0.105, precision
0.291, and recall 0.336 at the pixel level against the Bernsen pseudo-label. Considered alone, this
is a substantial drop from validation performance. However, the predicted total defect volume
fraction, 0.041\%, closely matches the pseudo-label's own volume fraction of 0.043\%
(Table~\ref{tab:podfam_classical}), while predicting fewer, larger, differently-distributed pore
instances (8,579 predicted pores vs.\ 11,928 labelled, mean diameter 6.35~px vs.\ 5.41~px). This
combination, low per-pixel overlap alongside near-identical total predicted volume, indicates the
detector has learned an approximately correct overall porosity level on data it never trained on,
while diverging from the pseudo-label's specific pore boundaries and instance count. Both the
pixel-level metrics and the volume-fraction agreement are reported together here, since either number
alone would misrepresent the result in a different direction.

\mysection{10.4 Unified Comparison}
Table~\ref{tab:podfam_unified} places the classical methods and the trained detector side by side on
the held-out test sample sample\_01.

\begin{table}[htbp]
\centering
\caption{Classical vs.\ U-Net, held-out test sample. U-Net's predicted defect fraction (0.041\%)
  sits between Otsu's near-zero result and Bernsen's label, closest to the label it was trained to
  reproduce, on data it never saw during training.}
\label{tab:podfam_unified}
\begin{tabular}{lllll}
\toprule
\textbf{Sample} & \textbf{Method} & \textbf{Defect (\%)} & \textbf{Pore count} & \textbf{Mean $\varnothing$ (px)} \\
\midrule
sample\_01 (held-out test) & Otsu             & 0.003  & 997       & 5.36 \\
                            & Yen              & 42.535 & 1,743,429 & 5.15 \\
                            & Bernsen (label)  & 0.043  & 11,928    & 5.41 \\
                            & U-Net (pred.)    & 0.041  & 8,579     & 6.35 \\
\bottomrule
\end{tabular}
\end{table}

\mysection{10.5 Extended Validation: sample\_07}
A sixth PODFAM sample, sample\_07 (1,473 slices, $1143\times1215$ pixels), was acquired and processed
after the five-sample validation reported above, using the identical preprocessing, thresholding, and
trained-detector inference pipeline described in Section~9, with no retraining or checkpoint changes.
sample\_07 was not assigned a role in the ranked train/validation/test split (Section~9.4); it is
reported here purely as an out-of-distribution check on the same epoch-8 checkpoint used throughout
Section~10.

\begin{table}[htbp]
\centering
\caption{sample\_07, full 1,473-slice volume. All three classical methods report 17--19\% defect
  fraction; the U-Net predicts 3.052\%.}
\label{tab:sample07_classical}
\begin{tabular}{llll}
\toprule
\textbf{Method} & \textbf{Defect (\%)} & \textbf{Pore count} & \textbf{Mean $\varnothing$ (px)} \\
\midrule
Otsu           & 18.177 & 330,344 & 9.15 \\
Yen            & 17.581 & 57,380  & 19.19 \\
Bernsen        & 19.369 & 354,426 & 9.15 \\
U-Net (pred.)  & 3.052  & 744,206 & 6.43 \\
\bottomrule
\end{tabular}
\end{table}

All three classical methods agree with one another on sample\_07 (17.6--19.4\%), which superficially
looks like corroboration rather than error. However, the same diagnostic applied in Section~9.5
identifies the same failure mode: sample\_07's processed slices have a mean intensity of 108.3
(measured on the mid-volume slice), below the fixed intensity threshold (128) that Bernsen's
low-contrast fallback branch uses to separate solid material from void. The auto-computed DCT for
sample\_07 (227) is, as with every other PODFAM sample, far above the approximately 15 the underlying
method reports as typical~\citep{Kim2017}, so the fallback branch dominates the classification.
sample\_05 (mean intensity 79.9, Section~9.5) showed the identical pattern at a more extreme value
and was excluded via the 10\%-defect-fraction guard; sample\_07's mean intensity is closer to the
fallback's 128 boundary, so the resulting over-estimate (17--19\%) is smaller in magnitude than
sample\_05's (approximately 52\%) but attributable to the same upstream cause.
Table~\ref{tab:sample07_diagnostic} places this in context against the four samples from
Section~10.1.

\begin{table}[htbp]
\centering
\caption{Mean processed-slice intensity vs.\ Bernsen defect fraction across all PODFAM samples. The
  two samples whose mean intensity falls below the fallback threshold (128) are exactly the two
  samples with implausibly high Bernsen defect fractions, reproducing the same failure mode at two
  different severities rather than two unrelated anomalies.}
\label{tab:sample07_diagnostic}
\begin{tabular}{lllll}
\toprule
\textbf{Sample} & \textbf{Mean intensity} & \textbf{Auto-DCT} & \textbf{Below 128?} & \textbf{Bernsen defect (\%)} \\
\midrule
sample\_01           & 167.0 & 84      & No  & 0.043 \\
sample\_02           & 142.3 & 329     & No  & 7.586 \\
sample\_04           & 163.4 & 152     & No  & 1.994 \\
sample\_07           & 108.3 & 227     & Yes & 19.369 \\
sample\_05 (excl.)   & 79.9  & 329--520 & Yes & 52.767 \\
\bottomrule
\end{tabular}
\end{table}

Under this reading, the U-Net's prediction of 3.052\% is the more credible estimate for sample\_07:
it is of the same order as sample\_02's Bernsen-labelled fraction (7.586\%) and sample\_04's U-Net
prediction (3.239\%), rather than an order of magnitude higher like the three classical methods on
this sample. This is, however, an inference-only result, sample\_07 was never assigned a
ground-truth role, so this comparison should be read as a plausibility check on the detector's
generalisation, not as a validated accuracy figure; see Section~12 for the corresponding limitation.
Figures~\ref{fig:s07_mesh}--\ref{fig:s07_mosaic} show the U-Net's predicted defect volume for
sample\_07 from four complementary viewpoints.

\begin{figure}[htbp]
\centering
\includegraphics[width=0.75\linewidth]{sample_07_unet_3d_mesh.png}
\caption{sample\_07, U-Net-predicted defect volume, marching-cubes 3D mesh reconstruction.}
\label{fig:s07_mesh}
\end{figure}

\begin{figure}[htbp]
\centering
\includegraphics[width=0.75\linewidth]{sample_07_unet_defect_scatter.png}
\caption{sample\_07, U-Net-predicted defect voxel scatter, coloured by slice depth.}
\label{fig:s07_scatter}
\end{figure}

\begin{figure}[htbp]
\centering
\includegraphics[width=0.95\linewidth]{sample_07_unet_orthographic_mip.png}
\caption{sample\_07, orthographic maximum-intensity projections (top row) and U-Net defect density
  projections (bottom row) along the three principal axes.}
\label{fig:s07_mip}
\end{figure}

\begin{figure}[htbp]
\centering
\includegraphics[width=0.95\linewidth]{sample_07_unet_slice_mosaic.png}
\caption{sample\_07, mosaic of 20 evenly-spaced slices with the U-Net's predicted defect mask
  overlaid in red.}
\label{fig:s07_mosaic}
\end{figure}
